\chapter*{Abstract}
\addcontentsline{toc}{chapter}{Abstract}
\begin{singlespace}

% This dissertation investigates optimization strategies for deploying Large Language Models (LLMs) in enterprise applications. Drawing on sixteen publications (eleven peer-reviewed, five under review) across established international AI and NLP venues, the research advances three interconnected pillars: Retrieval-Augmented Generation (RAG) optimization, agentic AI for workflow automation, and enterprise deployment strategies.

% The first pillar establishes RAG optimization through systematic model evaluation and novel metrics. Our evaluation demonstrates that smaller, efficient models achieve competitive performance: Orca-2-13b attains the highest overall score among evaluated models with approximately 99\% answer relevancy---surpassing GPT-3.5-Turbo and GPT-4 on faithfulness while maintaining 33 tokens per second on consumer-grade GPUs. Extending to Llama 2 across 500 MS MARCO queries, we confirm that open-source models deliver superior answer quality with 24--49\% inference improvements on newer GPU architectures. To address content repetition, we introduce Repetition-Aware Performance (RAP), enabling up to 93.1\% repetition reduction with only 3.7\% performance degradation across twelve LLMs (2B--70B parameters).

% The second pillar presents autonomous multi-agent systems for enterprise workflow automation. Our invoice reconciliation system achieves 99.90\% success rate with edge-deployed Qwen2.5-32B at only 523 Joules per task. A 12-category error taxonomy reveals tool initialization failures as the primary reliability bottleneck, with catastrophic rates in small models (89\% errors at 3B) absent at 32B scale. Our LLM-as-a-Judge (LAJ) framework introduces Evaluation Completion Rate (ECR@1), enabling systematic accuracy-reliability-cost trade-off analysis across 20 model configurations spanning a 175$\times$ cost range (\$0.45--\$78.96 per 1K evaluations); GPT-4o Mini emerges as production-optimal, achieving 78$\times$ cost reduction versus GPT-5 while maintaining superior accuracy. Building on this, our Cooperative Multi-Agent System for Gherkin Generation (CMAS4G2)---using Microsoft's AutoGen Swarm---shows open-weight models approach proprietary performance (87.6\% vs.\ 90.4\% success, 2.8 pp gap). Critically, reasoning configuration profoundly impacts coordination: structured reasoning enables transformative gains for smaller models (Qwen3-8B: 49.9\% $\rightarrow$ 77.5\%) while excessive depth causes catastrophic degradation. Our Hierarchical Multi-Agent System for Payments (HMASP)---the first LLM-based agentic payment framework---achieves 97.0\% task success coordinating card registration and 3D-Secure payment processing. EdgeAgent demonstrates semi-automatic annotation where qwen3:32b achieves 98\% Agentic Success Rate (ASR) matching GPT-4.1 quality.

% The third pillar examines deployment strategies validating cross-domain generalizability. Benchmarking across six hardware platforms, we establish that Windows Subsystem for Linux achieves up to 21$\times$ speedup for smaller models, while medium-sized models (7B--32B) achieve optimal performance-efficiency trade-offs with qwen2.5:32b attaining 96.75\% F1 score. A central finding from 504 configurations across seven model families: reasoning effectiveness is systematically task-dependent---binary sentiment classification suffers degradation (up to $-$19.9 F1 pp) while 27-class emotion recognition benefits (+16.0 pp), with reasoning justified only where accuracy gains outweigh 2.1--54$\times$ computational overhead. Cross-domain applications in machine translation (Qwen2-72B: 0.7564 COMET with LoRA/QLoRA fine-tuning) and sports betting (279.48\% betting return with LLM-derived sentiment features) validate technique generalizability.

% This research establishes that open-source LLMs, when properly optimized, achieve performance competitive with proprietary alternatives while enabling privacy-preserving edge deployment. The contributions---novel metrics (RAP, VCR, ECR@1, ASR), architectural patterns for multi-agent systems, and empirically validated deployment guidelines---provide a foundation for realizing the transformative potential of generative AI in enterprise applications.

This dissertation investigates how to deploy Large Language Models (LLMs) effectively in enterprise settings, where accuracy, reliability, cost, privacy, and operational constraints often matter more than benchmark performance alone. Drawing on seventeen peer-reviewed publications (eleven published and six accepted for publication), the work develops and validates optimization strategies across three connected themes: retrieval-augmented generation (RAG), agentic AI for workflow automation, and deployment guidelines for real-world enterprise environments.

First, we study RAG optimization through systematic evaluation of open and proprietary models, highlighting conditions under which efficient open-weight models can match or exceed proprietary alternatives. To address a pervasive failure mode in open-source LLMs---output repetition---we introduce \emph{Repetition-Aware Performance (RAP)}, a metric that integrates repetition penalties into task performance. Across models ranging from 2B to 70B parameters, RAP-guided tuning reduces repetition by up to 93.1\% with only 3.7\% performance degradation, enabling more dependable RAG outputs in production.

Second, we develop multi-agent architectures for enterprise automation and propose reliability-centric evaluation for tool-using systems. An invoice reconciliation system achieves 99.9\% task success with edge-deployed models at 523 Joules per task, while a structured diagnostic taxonomy identifies tool initialization as a dominant source of agent failures in smaller models. For scalable automated assessment, we introduce an LLM-as-a-Judge framework and the reliability metric \emph{Evaluation Completion Rate (ECR@1)}, enabling principled accuracy-reliability-cost trade-offs across a 175$\times$ pricing range; GPT-4o Mini emerges as production-optimal at 78$\times$ cost reduction versus premium models. A cooperative multi-agent system for acceptance test generation (CMAS4G2) demonstrates that open-weight models approach proprietary performance (87.6\% vs.\ 90.4\% success), with reasoning configuration profoundly impacting coordination. For payment processing, we design the first LLM-based agentic payment framework (HMASP), achieving 97.0\% task success, and introduce \emph{Agentic Success Rate (ASR)}---a trajectory-level workflow fidelity metric grounded in process-mining conformance checking. Across 18 LLMs and 90{,}000 task instances, ASR exposes systematic procedural deviations invisible to conventional metrics: most models bypass confirmation checkpoints during payment checkout despite near-perfect task success, with regulatory implications for PCI-DSS compliance. Ongoing work is evolving the research prototype into ConvPayMAS---an agentic payment system implementing Google's AP2 three-mandate cryptographic verification with patent-pending security features (demonstrated at AAMAS 2026)---upon which an enterprise agentic evaluation and optimization framework is being designed.

Third, we examine deployment strategies across heterogeneous hardware and tasks. Benchmarking across six hardware platforms shows that platform choices (e.g., WSL vs.\ native Windows) can yield up to 21$\times$ inference speedups, while medium-sized models (7B--32B) achieve optimal performance-efficiency trade-offs. A central finding from 504 configurations across seven model families is that reasoning effectiveness is strongly task-dependent: reasoning degrades simple binary sentiment classification (up to $-19.9$ F1 pp, 100\% failure rate) yet substantially improves complex emotion recognition (up to $+16.0$ F1 pp, only 14\% failure rate). Entity matching evaluation across 46 model configurations and 8 benchmark datasets reveals that reasoning effects are also family-dependent---low reasoning is optimal for GPT-5, high reasoning is critical for open-weight models, and Claude models exhibit extreme asymmetric sensitivity---while open-weight models rival proprietary APIs (GPT-OSS:120b outperforming GPT-4.1 at 20$\times$ lower cost). Cross-domain studies in machine translation, logical reasoning, and sports analytics further validate generalizability of the proposed metrics and deployment principles.

Overall, this dissertation demonstrates that open-source LLMs, when optimized and evaluated with reliability-aware metrics, can support privacy-preserving and cost-effective enterprise deployment. The proposed metrics, architectural patterns, and empirically grounded guidelines provide a practical foundation for building robust generative AI systems in production settings.

\vspace{0.8em}
\noindent\textbf{Keywords:} Large Language Models, Retrieval-Augmented Generation, Multi-Agent Systems, Agentic Commerce, Reasoning, Edge Deployment, Enterprise AI, Model Quantization, Few-Shot Learning, Fine-Tuning, Sentiment Analysis, Entity Matching, Tool Invocation, Evaluation Metrics, LoRA, Privacy-Preserving AI, Workflow Automation

\end{singlespace}